
\section{Data and branch definition}
\label{sec:data}

SPARC provides the rotation-curve data context and comprises 175 galaxies
\citep{Lelli2016}. The computational workflow uses two explicitly separated
samples: an $N=134$ partial-rank branch and an $N=76$ spectral-control branch.
The branch sizes are not interchangeable.

\subsection{Branch membership}

The frozen partial-rank branch required completion of the reconstruction
contract with at least eight valid radii, finite required vectors, positive
radial and velocity fields, a finite radial-template covariate, and a
computable radius-conditioned statistic. The statistic-level implementation
separately requires at least five finite joint points.

The spectral-control branch is the subset with at least 16 reconstructed
points, at least 16 finite unique positive-radius samples, finite interpolation
to a fixed 64-point log-radius grid, and nonzero finite spectral power.

The radius-conditioning coordinate is the dimensionless observed-radius
fraction
\begin{equation}
r_{gi}=\frac{R_{gi}}{R_{\max,g}},
\qquad
R_{\max,g}=\max_i R_{gi},
\end{equation}
where $R_{gi}$ is taken from the observed rotation-curve radius column after
nonfinite or nonpositive radius and velocity rows have been excluded, and
$R_{\max,g}$ is the maximum retained observed radius for galaxy $g$. The same
coordinate is used in the post-reconstruction validity mask and partial-rank
residualization. Correctly constructed values lie in $0<r_{gi}\leq1$. The
adapter also applies a permissive upper check $r_{gi}\leq1.1$; this check is a
post-construction implementation guard and does not redefine the
normalization or admit any additional correctly constructed observed-radius
fraction.
